\subsection{Prime ideals of height 1 in a Krull domain}

DEFINITION 4. Let A be an integral domain. A prime ideal p of A is said to be of height 1 if it is minimal among the non-zero prime ideals of A.

We shall also say that the ideal (0) is of height 0; a prime ideal of height \(\leqslant 1\) is therefore by definition equal to (0) or of height 1.

We shall define below, in a general way, the height of a prime ideal.

THEOREM 3. Let A be a Krull domain and p an integral ideal of A. For p to be the divisorial ideal corresponding to an extremal divisor, it is necessary and sufficient that p be a prime ideal of height 1.

If p is the divisorial ideal corresponding to an extremal divisor, we know (no. 4, Corollary 1 to Proposition 6) that p is prime and that \(A_{p}\) is a discrete valuation ring; as \(A_{p}\) has no prime ideals other than (0) and \(pA_{p}\) , (0) and p are the only prime ideals of A contained in p (Chapter II, §3, no. 1, Proposition 3); hence p is of height 1. Conversely, we shall show first that every prime ideal \(p \neq (0)\) of A contains a divisorial prime ideal q corresponding to an extremal divisor: for, as \(A_{p} \neq K\) , \(A_{p}\) is the intersection of a non-empty family ( \(A_{t}\) ) of essential valuation rings (no. 4, Proposition 6); each \(A_{t}\) is of the form \(A_{q_{t}}\) (no. 4, Corollary 1 to Proposition 6) and from \(A_{p} \subset A_{q_{t}}\) we deduce that \(q_{t} \subset p\) . Thus, if p is of height 1, then p = q, which shows that p is the divisorial ideal corresponding to an extremal divisor.

COROLLARY 1. In a Krull domain every non-zero prime ideal m contains a prime ideal of height 1. If m is not of height 1, then div m = 0 and A: m = A.

The first assertion has already been seen in the course of the proof of Theorem 3. If \(m\) is not of height 1 and \(p\) is a prime ideal of height 1 contained in \(m\) , then \(p \subset \tilde{m}\) and \(p \neq \tilde{m}\) ; as \(\text{div } p\) is extremal, necessarily \(\text{div } m = \text{div } \tilde{m} = 0\) ; hence \(\text{div}(A: m) = 0\) and, as \(A: m\) is divisorial (no. 1, Proposition 1), \(A: m = A\) .

COROLLARY 2. Let A be a Krull domain, K its field of fractions, v a valuation on K which is positive on A and p the set of \(x \in A\) such that \(v(x) > 0\) . If the prime ideal p is of height 1, v is equivalent to an essential valuation of A.

Let \(B\) be the ring of \(v\) and \(m\) its ideal. Then \(m \cap A = p\) and hence \(A_p \subset B\) . Now \(A_p\) is a discrete valuation ring (Theorem 3 and Corollary 1 to Proposition 6). As \(p \neq (0)\) , \(B \neq K\) and hence \(B = A_p\) (Chapter VI, § 4, no. 5, Proposition 6).

THEOREM 4. Let A be an integral domain and M the set of its prime ideals of height 1. For A to be a Krull domain, it is necessary and sufficient that the following properties are satisfied:

\begin{enumerate}
\def\labelenumi{(\roman{enumi})}
\item
  For all \(\mathfrak{p} \in \mathbf{M}\) , \(\mathbf{A}_{\mathfrak{p}}\) is a discrete valuation ring.
\item
  A is the intersection of the \(\mathbf{A}_{\mathfrak{p}}\) for \(\mathfrak{p} \in \mathbf{M}\) .
\item
  For all \(x \neq 0\) in \(\mathbf{A}\) , there exists only a finite number of ideals \(\mathfrak{p} \in \mathbf{M}\) such that \(x \in \mathfrak{p}\) .
\end{enumerate}

Moreover, the valuations corresponding to the \(A_{p}\) for \(p \in M\) are the essential valuations of A.

The conditions are trivially sufficient. Their necessity follows immediately from Theorem 3 of no. 4, Corollary 1 to Proposition 6 and the fact that the essential valuations of A satisfy the conditions of Definition 3 of no. 3.

PROPOSITION 10. Let A be an integrally closed Noetherian domain and a an integral ideal of A. The following conditions are equivalent:

\begin{enumerate}
\def\labelenumi{(\alph{enumi})}
\item
  a is divisorial;
\item
  the prime ideals associated with A/a are of height 1.
\end{enumerate}

Recall that, if \(a = \bigcap_{i=1}^{n} q_i\) is a reduced primary decomposition of a and \(p_i\) denotes the prime ideal corresponding to \(q_i\) , the prime ideals associated with A/a are just the \(p_i\) (Chapter IV, §2, no. 3, Proposition 4). The fact that (a) implies (b) then follows from Proposition 8 of no. 4. Conversely, if, in the above notation, the \(p_i\) are of height 1, \(A_{p_i}\) is a discrete valuation ring (Theorem 4); now, \(q_i = q_i A_{p_i} \cap A\) (Chapter IV, §2, no. 1, Proposition 3); denoting by \(v_i\) the essential valuation corresponding to \(p_i\) , there therefore exists an integer \(n_i\) such that \(q_i\) is the set of \(x \in A\) such that \(v_i(x) \geqslant n_i\) ; this shows that the \(q_i\) are divisorial (no. 4, Proposition 5), hence also is a.

